\IfFileExists{stacks-project.cls}{%
\documentclass{stacks-project}
}{%
\documentclass{amsart}
}

% For dealing with references we use the comment environment
\usepackage{verbatim}
\newenvironment{reference}{\comment}{\endcomment}
%\newenvironment{reference}{}{}
\newenvironment{slogan}{\comment}{\endcomment}
\newenvironment{history}{\comment}{\endcomment}

% For commutative diagrams we use Xy-pic
\usepackage[all]{xy}

% We use 2cell for 2-commutative diagrams.
\xyoption{2cell}
\UseAllTwocells

% We use multicol for the list of chapters between chapters
\usepackage{multicol}

% This is generally recommended for better output
\usepackage{lmodern}
\usepackage[T1]{fontenc}

% For cross-file-references
\usepackage{xr-hyper}

% Package for hypertext links:
\usepackage{hyperref}

% For any local file, say "hello.tex" you want to link to please
% use \externaldocument[hello-]{hello}
\externaldocument[introduction-]{introduction}
\externaldocument[conventions-]{conventions}
\externaldocument[sets-]{sets}
\externaldocument[categories-]{categories}
\externaldocument[topology-]{topology}
\externaldocument[sheaves-]{sheaves}
\externaldocument[sites-]{sites}
\externaldocument[stacks-]{stacks}
\externaldocument[fields-]{fields}
\externaldocument[algebra-]{algebra}
\externaldocument[brauer-]{brauer}
\externaldocument[homology-]{homology}
\externaldocument[derived-]{derived}
\externaldocument[simplicial-]{simplicial}
\externaldocument[more-algebra-]{more-algebra}
\externaldocument[smoothing-]{smoothing}
\externaldocument[modules-]{modules}
\externaldocument[sites-modules-]{sites-modules}
\externaldocument[injectives-]{injectives}
\externaldocument[cohomology-]{cohomology}
\externaldocument[sites-cohomology-]{sites-cohomology}
\externaldocument[dga-]{dga}
\externaldocument[dpa-]{dpa}
\externaldocument[sdga-]{sdga}
\externaldocument[hypercovering-]{hypercovering}
\externaldocument[schemes-]{schemes}
\externaldocument[constructions-]{constructions}
\externaldocument[properties-]{properties}
\externaldocument[morphisms-]{morphisms}
\externaldocument[coherent-]{coherent}
\externaldocument[divisors-]{divisors}
\externaldocument[limits-]{limits}
\externaldocument[varieties-]{varieties}
\externaldocument[topologies-]{topologies}
\externaldocument[descent-]{descent}
\externaldocument[perfect-]{perfect}
\externaldocument[more-morphisms-]{more-morphisms}
\externaldocument[flat-]{flat}
\externaldocument[groupoids-]{groupoids}
\externaldocument[more-groupoids-]{more-groupoids}
\externaldocument[etale-]{etale}
\externaldocument[chow-]{chow}
\externaldocument[intersection-]{intersection}
\externaldocument[pic-]{pic}
\externaldocument[weil-]{weil}
\externaldocument[adequate-]{adequate}
\externaldocument[dualizing-]{dualizing}
\externaldocument[duality-]{duality}
\externaldocument[discriminant-]{discriminant}
\externaldocument[derham-]{derham}
\externaldocument[local-cohomology-]{local-cohomology}
\externaldocument[algebraization-]{algebraization}
\externaldocument[curves-]{curves}
\externaldocument[resolve-]{resolve}
\externaldocument[models-]{models}
\externaldocument[functors-]{functors}
\externaldocument[equiv-]{equiv}
\externaldocument[pione-]{pione}
\externaldocument[etale-cohomology-]{etale-cohomology}
\externaldocument[proetale-]{proetale}
\externaldocument[relative-cycles-]{relative-cycles}
\externaldocument[more-etale-]{more-etale}
\externaldocument[trace-]{trace}
\externaldocument[crystalline-]{crystalline}
\externaldocument[spaces-]{spaces}
\externaldocument[spaces-properties-]{spaces-properties}
\externaldocument[spaces-morphisms-]{spaces-morphisms}
\externaldocument[decent-spaces-]{decent-spaces}
\externaldocument[spaces-cohomology-]{spaces-cohomology}
\externaldocument[spaces-limits-]{spaces-limits}
\externaldocument[spaces-divisors-]{spaces-divisors}
\externaldocument[spaces-over-fields-]{spaces-over-fields}
\externaldocument[spaces-topologies-]{spaces-topologies}
\externaldocument[spaces-descent-]{spaces-descent}
\externaldocument[spaces-perfect-]{spaces-perfect}
\externaldocument[spaces-more-morphisms-]{spaces-more-morphisms}
\externaldocument[spaces-flat-]{spaces-flat}
\externaldocument[spaces-groupoids-]{spaces-groupoids}
\externaldocument[spaces-more-groupoids-]{spaces-more-groupoids}
\externaldocument[bootstrap-]{bootstrap}
\externaldocument[spaces-pushouts-]{spaces-pushouts}
\externaldocument[spaces-chow-]{spaces-chow}
\externaldocument[groupoids-quotients-]{groupoids-quotients}
\externaldocument[spaces-more-cohomology-]{spaces-more-cohomology}
\externaldocument[spaces-simplicial-]{spaces-simplicial}
\externaldocument[spaces-duality-]{spaces-duality}
\externaldocument[formal-spaces-]{formal-spaces}
\externaldocument[restricted-]{restricted}
\externaldocument[spaces-resolve-]{spaces-resolve}
\externaldocument[formal-defos-]{formal-defos}
\externaldocument[defos-]{defos}
\externaldocument[cotangent-]{cotangent}
\externaldocument[examples-defos-]{examples-defos}
\externaldocument[algebraic-]{algebraic}
\externaldocument[examples-stacks-]{examples-stacks}
\externaldocument[stacks-sheaves-]{stacks-sheaves}
\externaldocument[criteria-]{criteria}
\externaldocument[artin-]{artin}
\externaldocument[quot-]{quot}
\externaldocument[stacks-properties-]{stacks-properties}
\externaldocument[stacks-morphisms-]{stacks-morphisms}
\externaldocument[stacks-limits-]{stacks-limits}
\externaldocument[stacks-cohomology-]{stacks-cohomology}
\externaldocument[stacks-perfect-]{stacks-perfect}
\externaldocument[stacks-introduction-]{stacks-introduction}
\externaldocument[stacks-more-morphisms-]{stacks-more-morphisms}
\externaldocument[stacks-geometry-]{stacks-geometry}
\externaldocument[moduli-]{moduli}
\externaldocument[moduli-curves-]{moduli-curves}
\externaldocument[examples-]{examples}
\externaldocument[exercises-]{exercises}
\externaldocument[guide-]{guide}
\externaldocument[desirables-]{desirables}
\externaldocument[coding-]{coding}
\externaldocument[obsolete-]{obsolete}
\externaldocument[fdl-]{fdl}
\externaldocument[index-]{index}

% Theorem environments.
%
\theoremstyle{plain}
\newtheorem{theorem}[subsection]{Theorem}
\newtheorem{proposition}[subsection]{Proposition}
\newtheorem{lemma}[subsection]{Lemma}

\theoremstyle{definition}
\newtheorem{definition}[subsection]{Definition}
\newtheorem{example}[subsection]{Example}
\newtheorem{exercise}[subsection]{Exercise}
\newtheorem{situation}[subsection]{Situation}

\theoremstyle{remark}
\newtheorem{remark}[subsection]{Remark}
\newtheorem{remarks}[subsection]{Remarks}

\numberwithin{equation}{subsection}

% Macros
%
\def\lim{\mathop{\mathrm{lim}}\nolimits}
\def\colim{\mathop{\mathrm{colim}}\nolimits}
\def\Spec{\mathop{\mathrm{Spec}}}
\def\Hom{\mathop{\mathrm{Hom}}\nolimits}
\def\Ext{\mathop{\mathrm{Ext}}\nolimits}
\def\SheafHom{\mathop{\mathcal{H}\!\mathit{om}}\nolimits}
\def\SheafExt{\mathop{\mathcal{E}\!\mathit{xt}}\nolimits}
\def\Sch{\mathit{Sch}}
\def\Mor{\mathop{\mathrm{Mor}}\nolimits}
\def\Ob{\mathop{\mathrm{Ob}}\nolimits}
\def\Sh{\mathop{\mathit{Sh}}\nolimits}
\def\NL{\mathop{N\!L}\nolimits}
\def\CH{\mathop{\mathrm{CH}}\nolimits}
\def\proetale{{pro\text{-}\acute{e}tale}}
\def\etale{{\acute{e}tale}}
\def\QCoh{\mathit{QCoh}}
\def\Ker{\mathop{\mathrm{Ker}}}
\def\Im{\mathop{\mathrm{Im}}}
\def\Coker{\mathop{\mathrm{Coker}}}
\def\Coim{\mathop{\mathrm{Coim}}}

% Boxtimes
%
\DeclareMathSymbol{\boxtimes}{\mathbin}{AMSa}{"02}

%
% Macros for moduli stacks/spaces
%
\def\QCohstack{\mathcal{QC}\!\mathit{oh}}
\def\Cohstack{\mathcal{C}\!\mathit{oh}}
\def\Spacesstack{\mathcal{S}\!\mathit{paces}}
\def\Quotfunctor{\mathrm{Quot}}
\def\Hilbfunctor{\mathrm{Hilb}}
\def\Curvesstack{\mathcal{C}\!\mathit{urves}}
\def\Polarizedstack{\mathcal{P}\!\mathit{olarized}}
\def\Complexesstack{\mathcal{C}\!\mathit{omplexes}}
% \Pic is the operator that assigns to X its picard group, usage \Pic(X)
% \Picardstack_{X/B} denotes the Picard stack of X over B
% \Picardfunctor_{X/B} denotes the Picard functor of X over B
\def\Pic{\mathop{\mathrm{Pic}}\nolimits}
\def\Picardstack{\mathcal{P}\!\mathit{ic}}
\def\Picardfunctor{\mathrm{Pic}}
\def\Deformationcategory{\mathcal{D}\!\mathit{ef}}
\usepackage{cleveref}
\crefname{originalreference}{original reference}{original references}
\Crefname{originalreference}{Original reference}{Original references}
\makeatletter
\newlabel{algebra-lemma-NAK}{{00DV}{source}{}{sourcepointer-1}{}}
\newlabel{algebra-lemma-NAK@cref}{{[originalreference][0][]00DV}{[1][0][]source}{}{}{}}
\makeatother
\crefname{originalreference}{original reference}{original references}
\Crefname{originalreference}{Original reference}{Original references}
\makeatletter
\newlabel{algebra-lemma-localize-NL}{{00S7}{source}{}{sourcepointer-2}{}}
\newlabel{algebra-lemma-localize-NL@cref}{{[originalreference][0][]00S7}{[1][0][]source}{}{}{}}
\makeatother
\crefname{originallemma}{original lemma}{original lemmas}
\Crefname{originallemma}{Original lemma}{Original lemmas}
\makeatletter
\newlabel{lemma-NL-is-completion}{{0GAF}{source}{}{sourcepointer-3}{}}
\newlabel{lemma-NL-is-completion@cref}{{[originallemma][0][]0GAF}{[1][0][]source}{}{}{}}
\makeatother
\crefname{originallemma}{original lemma}{original lemmas}
\Crefname{originallemma}{Original lemma}{Original lemmas}
\makeatletter
\newlabel{lemma-equivalent-with-artin-smooth}{{0GAJ}{source}{}{sourcepointer-4}{}}
\newlabel{lemma-equivalent-with-artin-smooth@cref}{{[originallemma][0][]0GAJ}{[1][0][]source}{}{}{}}
\makeatother
\crefname{originalreference}{original reference}{original references}
\Crefname{originalreference}{Original reference}{Original references}
\makeatletter
\newlabel{more-algebra-lemma-base-change-RHom}{{0A6A}{source}{}{sourcepointer-5}{}}
\newlabel{more-algebra-lemma-base-change-RHom@cref}{{[originalreference][0][]0A6A}{[1][0][]source}{}{}{}}
\makeatother
\crefname{originalreference}{original reference}{original references}
\Crefname{originalreference}{Original reference}{Original references}
\makeatletter
\newlabel{more-algebra-lemma-base-change-property-ext-1-annihilated}{{0G9H}{source}{}{sourcepointer-6}{}}
\newlabel{more-algebra-lemma-base-change-property-ext-1-annihilated@cref}{{[originalreference][0][]0G9H}{[1][0][]source}{}{}{}}
\makeatother
\crefname{originalreference}{original reference}{original references}
\Crefname{originalreference}{Original reference}{Original references}
\makeatletter
\newlabel{more-algebra-lemma-ext-1-annihilated}{{0G9K}{source}{}{sourcepointer-7}{}}
\newlabel{more-algebra-lemma-ext-1-annihilated@cref}{{[originalreference][0][]0G9K}{[1][0][]source}{}{}{}}
\makeatother
\crefname{originalreference}{original reference}{original references}
\Crefname{originalreference}{Original reference}{Original references}
\makeatletter
\newlabel{more-algebra-lemma-pseudo-coherence-and-base-change-ext}{{087R}{source}{}{sourcepointer-8}{}}
\newlabel{more-algebra-lemma-pseudo-coherence-and-base-change-ext@cref}{{[originalreference][0][]087R}{[1][0][]source}{}{}{}}
\makeatother
\title{Smoothness off V(I) versus rig-smoothness of the completion}
\author{The Stacks Project Authors}
\date{stacks-project:restricted.tex:local-snapshot-20261012; GNU FDL 1.2 or later}
\begin{document}
\maketitle
\noindent\textit{Editorial note.} Counted span is the author's own proof environment, source lines 955-1013, following the statement at 942-953.\par\medskip
\noindent\textit{Editorial note.} Both 'small detail omitted' flags are readings through lemma-equivalent-with-artin-smooth, whose use 'without further mention' the author declares at the start; flagged in the notes.\par\medskip
\noindent\textit{Editorial note.} Steps the author leaves to the reader are flagged here for the reviewer: source line(s) 967, 1007 carry an "omitted"/"follows easily"/"immediate from the definitions" phrase. Those sentences are retained verbatim; the displayed argument around them stays inside the counted span.\par\medskip
\noindent\textit{Allowed standard prerequisites (editorial).}\begin{itemize}
\item Equivalent conditions for rig-smoothness. lemma-equivalent-with-artin-smooth (same chapter, used without further mention; its proof has stated omissions and was not counted). Translates between annihilated Ext classes and the dotted-arrow condition.
\item Completion of the naive cotangent complex. lemma-NL-is-completion (same chapter, cited at \textasciitilde{}965) and More on Algebra lemmas base-change-property-ext-1-annihilated, base-change-RHom (cited at \textasciitilde{}966-982). Relates NL of B and of the completion.
\item Nakayama and Krull tools. Algebra, Lemma algebra-lemma-NAK (cited at \textasciitilde{}991). Produces the element g = 1 + x annihilating the class module.
\end{itemize}\medskip
\ifdefined\hypertarget\hypertarget{sourcepointer-1}{}\fi
\noindent\textit{Source reference.} Original reference 00DV: focused-seeds/geometry/S9-restricted-rig-smooth/ref-tags.txt, source lines 5--5. Navigation pointer only: the counted proof cites this statement by \textbackslash{}ref as a standing input; the deduction the counted proof has to make is carried out inside the counted span, and no part of that deduction is delegated to this pointer.\par\medskip
\ifdefined\hypertarget\hypertarget{sourcepointer-2}{}\fi
\noindent\textit{Source reference.} Original reference 00S7: focused-seeds/geometry/S9-restricted-rig-smooth/ref-tags.txt, source lines 6--6. Navigation pointer only: the counted proof cites this statement by \textbackslash{}ref as a standing input; the deduction the counted proof has to make is carried out inside the counted span, and no part of that deduction is delegated to this pointer.\par\medskip
\ifdefined\hypertarget\hypertarget{sourcepointer-3}{}\fi
\noindent\textit{Source reference.} Original lemma 0GAF: sources/stacks-restricted/pinned-original-source.tex, source lines 560--560. Navigation pointer only: the counted proof cites this statement by \textbackslash{}ref as a standing input; the deduction the counted proof has to make is carried out inside the counted span, and no part of that deduction is delegated to this pointer.\par\medskip
\ifdefined\hypertarget\hypertarget{sourcepointer-4}{}\fi
\noindent\textit{Source reference.} Original lemma 0GAJ: sources/stacks-restricted/pinned-original-source.tex, source lines 843--843. Navigation pointer only: the counted proof cites this statement by \textbackslash{}ref as a standing input; the deduction the counted proof has to make is carried out inside the counted span, and no part of that deduction is delegated to this pointer.\par\medskip
\ifdefined\hypertarget\hypertarget{sourcepointer-5}{}\fi
\noindent\textit{Source reference.} Original reference 0A6A: focused-seeds/geometry/S9-restricted-rig-smooth/ref-tags.txt, source lines 7--7. Navigation pointer only: the counted proof cites this statement by \textbackslash{}ref as a standing input; the deduction the counted proof has to make is carried out inside the counted span, and no part of that deduction is delegated to this pointer.\par\medskip
\ifdefined\hypertarget\hypertarget{sourcepointer-6}{}\fi
\noindent\textit{Source reference.} Original reference 0G9H: focused-seeds/geometry/S9-restricted-rig-smooth/ref-tags.txt, source lines 8--8. Navigation pointer only: the counted proof cites this statement by \textbackslash{}ref as a standing input; the deduction the counted proof has to make is carried out inside the counted span, and no part of that deduction is delegated to this pointer.\par\medskip
\ifdefined\hypertarget\hypertarget{sourcepointer-7}{}\fi
\noindent\textit{Source reference.} Original reference 0G9K: focused-seeds/geometry/S9-restricted-rig-smooth/ref-tags.txt, source lines 9--9. Navigation pointer only: the counted proof cites this statement by \textbackslash{}ref as a standing input; the deduction the counted proof has to make is carried out inside the counted span, and no part of that deduction is delegated to this pointer.\par\medskip
\ifdefined\hypertarget\hypertarget{sourcepointer-8}{}\fi
\noindent\textit{Source reference.} Original reference 087R: focused-seeds/geometry/S9-restricted-rig-smooth/ref-tags.txt, source lines 10--10. Navigation pointer only: the counted proof cites this statement by \textbackslash{}ref as a standing input; the deduction the counted proof has to make is carried out inside the counted span, and no part of that deduction is delegated to this pointer.\par\medskip
\begin{lemma}
\label{lemma-rig-smooth}
Let $A$ be a Noetherian ring and let $I$ be an ideal.
Let $B$ be a finite type $A$-algebra.
\begin{enumerate}
\item If $\Spec(B) \to \Spec(A)$ is smooth over $\Spec(A) \setminus V(I)$,
then $B^\wedge$ is rig-smooth over $(A, I)$.
\item If $B^\wedge$ is rig-smooth over $(A, I)$,
then there exists $g \in 1 + IB$ such that $\Spec(B_g)$ is smooth
over $\Spec(A) \setminus V(I)$.
\end{enumerate}
\end{lemma}

\begin{proof}
We will use Lemma \ref{lemma-equivalent-with-artin-smooth}
without further mention.

\medskip\noindent
Assume (1). Recall that formation of $\NL_{B/A}$ commutes with
localization, see Algebra, Lemma \ref{algebra-lemma-localize-NL}.
Hence by the very definition of smooth ring maps (in terms of
the naive cotangent complex being quasi-isomorphic to a finite
projective module placed in degree $0$), we see that
$\NL_{B/A}$ satisfies the fourth equivalent condition
of More on Algebra, Lemma \ref{more-algebra-lemma-ext-1-annihilated}
with respect to the ideal $IB$ (small detail omitted).
Since $\NL_{B^\wedge/A}^\wedge = \NL_{B/A} \otimes_B B^\wedge$
by Lemma \ref{lemma-NL-is-completion} we conclude (2) holds by
More on Algebra, Lemma
\ref{more-algebra-lemma-base-change-property-ext-1-annihilated}.

\medskip\noindent
Assume (2). Choose a presentation
$B = A[x_1, \ldots, x_n]/J$, set $N = J/J^2$, and
consider the element $\xi \in \Ext^1_B(\NL_{B/A}, J/J^2)$
determined by the identity map on $J/J^2$.
Using again that $\NL_{B^\wedge/A}^\wedge = \NL_{B/A} \otimes_B B^\wedge$
we find that our assumption implies the image
$$
\xi \otimes 1 \in
\Ext^1_{B^\wedge}(\NL_{B/A} \otimes_B B^\wedge, N \otimes_B B^\wedge) =
\Ext^1_{B^\wedge}(\NL_{B/A}, N) \otimes_B B^\wedge
$$
is annihilated by $I^c$ for some integer $c \geq 0$.
The equality holds for example by
More on Algebra, Lemma \ref{more-algebra-lemma-base-change-RHom}
(but can also easily be deduced from the much simpler
More on Algebra, Lemma
\ref{more-algebra-lemma-pseudo-coherence-and-base-change-ext}).
Thus $M = I^cB\xi \subset \Ext^1_B(\NL_{B/A}, N)$ is a finite submodule
which maps to zero in $\Ext^1_B(\NL_{B/A}, N) \otimes_B B^\wedge$.
Since $B \to B^\wedge$ is flat this means that
$M \otimes_B B^\wedge$ is zero. By 
Nakayama's lemma (Algebra, Lemma \ref{algebra-lemma-NAK})
this means that $M = I^cB\xi$ is annihilated by an element
of the form $g = 1 + x$ with $x \in IB$.
This implies that for every $b \in I^cB$ there is a
$B$-linear dotted arrow making the diagram commute
$$
\xymatrix{
J/J^2 \ar[r] \ar[d]^b & \bigoplus B\text{d}x_i \ar@{..>}[d]^h \\
J/J^2 \ar[r] & (J/J^2)_g
}
$$
Thus $(\NL_{B/A})_{gb}$ is quasi-isomorphic
to a finite projective module; small detail omitted.
Since $(\NL_{B/A})_{gb} = \NL_{B_{gb}/A}$ in $D(B_{gb})$
this shows that $B_{gb}$ is smooth over $\Spec(A)$.
As this holds for all $b \in I^cB$ we conclude
that $\Spec(B_g) \to \Spec(A)$ is smooth over $\Spec(A) \setminus V(I)$
as desired.
\end{proof}

\end{document}
